\documentclass[compress, english, aspectratio=169]{beamer}

\usepackage{pythontex} % To run Python code
\usepackage[utf8]{inputenc}
\usepackage{amssymb}
%\usepackage{graphicx}
\graphicspath{{images/}}

\usepackage[scheme=plain]{ctex} % To display Chinese characters

\usetheme[navigation]{UMONS}
%\usetheme[navigation, no-subsection, no-totalframenumber]{UMONS}

\newcommand{\IR}{\mathbb{R}}

%% Presentation information
\title{Industrial Instrumentation: Design of Experiments}
\subtitle{GUI project with Python and PyQt5}
\author[V. \textsc{Stragier}]{Vincent \textsc{Stragier}}
\date{16th of June}
\institute[| FPMs]{%
  Polytechnic Faculty of Mons\\
  University of  Mons
  \\[2ex]
  \includegraphics[height=6ex]{LOGO_UMONS_EN}\hspace{2em}
  \includegraphics[height=6ex]{LOGO_FPMs}
}

\begin{document}
% Setup pythontex to be able to import another module located in other path
\begin{pyconcode}
import os;
os.chdir('../')
import sys
sys.path.append("..\\src")
sys.path.append(os.getcwd())
sys.path = list(set(sys.path))
\end{pyconcode}

%% Title
\begin{frame}[plain]
  \titlepage
\end{frame}

%% Table of content
\begin{frame}
  \tableofcontents
\end{frame}

%% Proverb
\begin{frame}
\frametitle{Chinese proverb}
\begin{block}{读万卷书不如行万里路。(Dú wànjuànshū bù rú xíng wànlǐlù)}
It's better to walk thousands of miles than to read thousands of books.
\end{block}
\end{frame}

%% Theory
\section{Theory about the 2 factorial Design of Experiments}
\subsection{Factorial design}
\begin{frame}
\frametitle{Factorial design}
\begin{block}{Goals}
\begin{itemize}
\item efficient measurement:
\begin{itemize}
\item quickly arrive at the best possible results
\item omit unnecessary trials
\item give results with the lowest uncertainty
\item progress without failure
\item establish a model for the studied phenomenon
\item discover the optimal solution
\end{itemize}
\end{itemize}
\end{block}

\begin{equation}
m = m_0^{n}
\end{equation}

where $m$ is the number of measurements, $m_0$ is the number of levels and $n$ is the number of factors that you want to adjust. For the program $m_0=2$ is assumed. It will define the number of measurements to take in the experimental space.
\end{frame}

\section{Python implementation}
\subsection{How to program}
\begin{frame}
\frametitle{How to program}
\begin{block}{Setup}
\begin{itemize}
\item sheets of paper and pen;
\item book(s) or internet connexion;
\item a computer;
\item an IDE.
\end{itemize}
\end{block}
\end{frame}

\begin{frame}
\frametitle{How to program}
\begin{block}{Methodology}
\begin{itemize}
\item up-to-date state-of-the-art;
\item scheme of your program;
\item use the classical naming standards;
\item testing and debugging method;
\item document your code;
\item implement the minimal functions first;
\item do your research before asking anything.
\end{itemize}
\end{block}
\end{frame}



\subsection{Core functions}
\subsubsection{Algorithm}
\begin{frame}
\frametitle{Core functions}
\framesubtitle{Algorithm}
\begin{align*}
\hat{\underline{a}} = \left(X^T X\right)^{-1} X^T \underline{y}
\end{align*}

\noindent where: 

\begin{description}
\item[$\hat{\underline{a}}$] is a vector of coefficients (of shape $p\times 1$)
\item[$\underline{y}$] is a vector of measures (of shape $p\times 1$)
\item[$X^T$] is the transposed form of the matrix $T$
\item[$\left(X^T X\right)$] is the square information matrix (of shape $p\times p$)
\item[$\left(X^T X\right)^{-1}$] is the square dispersion matrix (of shape $p\times p$)
\end{description}
\end{frame}

\subsubsection{Aim}
\begin{frame}
\frametitle{Core functions}
\framesubtitle{Aim}
The aim is to compute the estimation of the coefficients to provide a model for the experiment. In our case the following model can be used:

\begin{align*}
y = \hat{a}_0 + \sum^n_{i=1} \hat{a}_i x_i + \sum^n_{i=1} \sum^n_{j=1} \hat{a}_{ij} x_i {x}_j + ...
\end{align*}

To ease the resolution of the coefficients, $\hat{X}$ is defined as followed:
\begin{equation*}
\hat{X} = \left(X^T X\right)^{-1} X^T
\end{equation*}
\end{frame}

\subsubsection{Console full example}
\begin{frame}[fragile]
\frametitle{Core functions}
\framesubtitle{Console full example}
\begin{pyconsole}
import design_of_experiments as doe
X_hat = doe.gen_X_hat(n=3)
X_hat
\end{pyconsole}
\end{frame}

\begin{frame}[fragile]
\frametitle{Core functions}
\framesubtitle{Console full example (continued)}
\begin{pyconsole}
# Speed [mg/min]
y1 = [53, 122, 20, 125, 48, 70, 68, 134]
# Cobalt Content [ppm]
y2 = [4100, 3510, 3950, 1270, 4870, 2810, 7750, 3580]
# Compute the coefficients
import numpy as np
a_hat_1 = np.dot(X_hat, y1)
a_hat_2 = np.dot(X_hat, y2)
a_hat_1
a_hat_2
\end{pyconsole}
\end{frame}

\subsection{GUI programming}
\begin{frame}
\frametitle{GUI programming}
\begin{block}{Main steps}
\begin{itemize}
\item design the GUI on the Qt Designer tool
\item convert the .ui file generated in a Python script
\item create a new script that will import the generated script and make the links between the main script and the GUI script
\end{itemize}
\end{block}
\end{frame}

\section{Demonstration}
\subsection{Requirements}
\begin{frame}
\frametitle{Requirements}
\begin{block}{Program and modules}
\begin{itemize}
\item Python 3.8.0 or above
\item the source code (\url{https://github.com/2010019970909/design_of_experiments})
\item any working LaTeX installation
\item Python modules:\\
\begin{description}
\item[scikit] for the statistics
\item[numpy] for the math (matrix dot product, etc.)
\item[matplotlib] for the plots
\item[PyQt5] for the GUI
\item[PyQt5-tools] to further development
\end{description}
\end{itemize}
\end{block}

\pyb{py -3.8 -m pip install --upgrade sklearn numpy matplotlib PyQt5 PyQt5-tools}
\end{frame}


\end{document}
%%% Local Variables: 
%%% mode: latex
%%% TeX-master: t
%%% End: 
